\documentclass[11pt]{article}
\usepackage[T1]{fontenc}
\usepackage[utf8]{inputenc}
\usepackage{lmodern,amsmath,amssymb,amsthm,mathtools}
\usepackage[margin=1.1in]{geometry}
\usepackage{microtype}
\usepackage[colorlinks=true,linkcolor=blue,citecolor=blue,urlcolor=blue]{hyperref}
\newtheorem{theorem}{Theorem}[section]
\newtheorem{proposition}[theorem]{Proposition}
\newtheorem{lemma}[theorem]{Lemma}
\newtheorem{remark}[theorem]{Remark}
\numberwithin{equation}{section}
\newcommand{\R}{\mathbb R}
\newcommand{\Sn}{S^{n-1}}
\newcommand{\Lip}{\operatorname{Lip}}
\newcommand{\dist}{\operatorname{dist}}
\newcommand{\supp}{\operatorname{supp}}
\newcommand{\PV}{\operatorname{PV}}
\newcommand{\osc}{\operatorname{osc}}
\newcommand{\T}{\mathcal T}
\newcommand{\Hh}{\mathbb H}
\newcommand{\Bfun}{\mathrm B}
\DeclareMathOperator*{\esssup}{ess\,sup}
\setlength{\emergencystretch}{2em}

\title{Optimal obstacle regularity for stable operators\\
with minimally smooth obstacles}
\author{Kaiwen Jin\thanks{Contact: \href{mailto:yui.ui.siki@gmail.com}{\texttt{yui.ui.siki@gmail.com}}.}}
\date{12 November 2026}
\hypersetup{pdftitle={Optimal obstacle regularity for stable operators with minimally smooth obstacles},pdfauthor={Kaiwen Jin},pdfsubject={Regularity for nonlocal obstacle problems}}
\begin{document}
\maketitle
\begin{abstract}
We prove the interior $C^{1+s}$ estimate for obstacle problems driven by
symmetric stable operators with positive angular Lipschitz density, for every
obstacle in $C^\beta$ with noninteger $\beta>1+s$. The exterior values of the
solution are arbitrary bounded data. This gives the minimally smooth obstacle
threshold in the local anisotropic problem. After subtracting and localizing
the obstacle, the forcing need only be H\"older continuous. We compare the solution
with shifted half-space profiles in a known positive half-ball. A thin-strip
Dirichlet decomposition controls the constant forcing defect; a flat boundary
quotient estimate gives the outer variation. Blow-ups at actual free-boundary
points identify the quotient's signed trace, yielding improvement of the value bounds
and of the compensated tail. A scale dichotomy and an interior covering
argument give the full $C^{1+s}$ norm estimate.
\end{abstract}

\section{Introduction}\label{sec:intro}

The nonlocal obstacle problem asks for a function that lies above a prescribed
obstacle and is harmonic for a nonlocal elliptic operator wherever it separates
from that obstacle. It is also the stationary optimal stopping equation for a
stable jump process, with the obstacle playing the role of its payoff; see
\cite[Chapter 4]{FRRO}. The set where the two functions meet is unknown.
Regularity across its boundary is therefore stronger than an interior
estimate in the harmonic region. The characteristic detachment profile has
degree $1+s$ for an operator of order $2s$, making $C^{1+s}$ the natural
regularity to seek.

For the fractional Laplacian, Caffarelli, De Silva and Savin proved this
regularity with obstacles in $C^\beta$ for $\beta>1+s$
\cite[Theorem 6.1]{CDS}. Their argument uses a truncated frequency formula
for the extension equation. For general stable operators, the local result
of Ros-Oton, Torres-Latorre and Weidner gives every subcritical estimate
$C^q$, $q<1+s$, and the endpoint $C^{1+s}$ estimate when
$\beta>1+2s$ \cite[Corollary 1.3]{RTW}. Their estimates allow arbitrary
bounded exterior values. The remaining minimally smooth anisotropic question
is recorded in \cite[Open Question 4.1, Section 4.6.1]{FRRO}.

We prove the same local endpoint estimate for every $\beta>1+s$ under
angular Lipschitz control. More precisely, let
\begin{equation}\label{eq:intro-operator}
 Lu(x)=\PV\int_{\R^n}(u(x)-u(x+y))
                   \frac{a(y/|y|)}{|y|^{n+2s}}\,dy.
\end{equation}

\begin{theorem}\label{thm:main}
Let $n\geq2$, $0<s<1$, and let $\beta>1+s$ be noninteger. Suppose
$a\in\Lip(S^{n-1})$ is even and $0<\lambda\leq a\leq\Lambda$.
For every $\phi\in C^\beta(B_1)$ and every bounded continuous viscosity
solution of
\[
                         \min\{Lu,u-\phi\}=0\quad\hbox{in }B_1,
\]
one has
\begin{equation}\label{eq:main-estimate}
 \|u\|_{C^{1+s}(B_{1/2})}
       \leq C\left(\|\phi\|_{C^\beta(B_1)}+
                                  \|u\|_{L^\infty(\R^n)}\right).
\end{equation}
Here $C$ depends only on $n,s,\beta,\lambda,\Lambda$ and
$\|a\|_{\Lip(S^{n-1})}$.
\end{theorem}

Thus the obstacle, solution concept, operator sign and quantitative norm
are those of the local question. The estimate includes arbitrary bounded
exterior values and does not require the coincidence set to stay away from
$\partial B_1$. Its new range relative to the cited anisotropic endpoint
statement is $1+s<\beta\leq1+2s$.

The difficulty is visible after subtracting the obstacle. If its usable
smoothness is $1+s+\alpha$, the new forcing is only
$C^{1-s+\alpha}$, below Lipschitz regularity when $\alpha<s$.
The convexity and classification of blow-ups from \cite[Theorem 1.5]{RTW}
still provide a coarse half-space profile. Turning this coarse information
into an endpoint bound requires finer control of the contact layer. Merely
locating the free boundary does not order the values above it.

We obtain that control from a value sandwich between two shifted profiles.
It supplies a known positive half-ball above the thin uncertain layer. A
Dirichlet decomposition on this fixed flat geometry separates the layer's
harmonic influence from a part that vanishes on the flat boundary. Torsion
at the layer scale bounds the unknown constant forcing defect. The quotient
of the zero-exterior part by the distance to the flat boundary to the power
$s$ then has a first-order expansion. We prove the required flat quotient
estimate using only angular Lipschitz regularity and the one-dimensional
higher Liouville theorem of \cite[Proposition 3.11]{AR}.

The sign of the expansion is fixed by blowing up at actual free-boundary
points inside the layer. Matching these inner limits to the outer quotient
ties its trace to the true contact displacement. This also excludes positive
pockets below the limiting contact plane. The resulting value comparison and
tail improvement iterate. A first-failure choice of scale treats the
nondegenerate branch, while the complementary branch already has stronger
decay. Finally, derivatives of a localized solution satisfy linear equations
in the positive phase, which converts the one-point estimate into
\eqref{eq:main-estimate}.

The strict inequality $\beta>1+s$ supplies a positive H\"older margin in
this argument. We do not address the equality case or nonhomogeneous kernels.
Section~\ref{sec:linear} proves the flat boundary estimate;
Section~\ref{sec:flatness} proves improvement of flatness; and
Section~\ref{sec:interior} completes the localization, scale selection and
norm estimate. We use the positive operator convention
\eqref{eq:intro-operator}; the norms and weighted tails are specified at
their first use below.

\section{Linear estimates in a flat half-space}\label{sec:linear}

Throughout the paper the operator has the positive sign convention
\begin{equation}\label{eq:operator}
 Lu(x)=\PV\int_{\R^n}(u(x)-u(x+y))K(y)\,dy,
 \qquad K(y)=\frac{a(y/|y|)}{|y|^{n+2s}},
\end{equation}
where $a$ is even, $\lambda\leq a\leq\Lambda$, and
$\|a\|_{\Lip(\Sn)}\leq M$. Constants may depend on $n,s,\lambda,\Lambda,M$
and on the fixed auxiliary exponents. They do not depend on an individual
solution's bound at infinity unless that bound is displayed.
We write $x=(x',t)$ and $\Hh=\{t>0\}$.

For noninteger $\nu=k+\theta$, $k\geq0$ and $0<\theta<1$, $C^\nu$
has its usual inhomogeneous norm: the supremum norms of derivatives through
order $k$, together with the $\theta$-H\"older seminorm of the derivatives of
order $k$. Brackets $[\cdot]_{C^\nu}$ denote the homogeneous seminorm.
The weighted value norm used below is
\[
 \|w\|_{L^\infty_\mu}=\sup_{x\in\R^n}
                 \frac{|w(x)|}{(1+|x|)^\mu}.
\]
All equations in this section can be read distributionally. At the regularity
greater than $2s$ used in the nonlinear argument, this agrees with the
classical integral and the corresponding viscosity equation.

\subsection{Growing profiles and local estimates}

The functions $t_+^{1+s}$ and $x_i t_+^s$ are useful comparison functions,
but their full integrals in \eqref{eq:operator} need not converge at infinity.
We use their operator only up to an additive constant. More precisely, for a
smooth radial cutoff $\chi_R$, equal to one on $B_R$ and zero outside $B_{2R}$,
the assertion $LU=F$ modulo constants means that there are constants $c_R$
such that
\[
 L(\chi_R U)-c_R\longrightarrow F
\]
on compact subsets of the domain, in distributions. The estimates below
also give convergence of the oscillatory far part. This is the $k=1$
interpretation of \cite[Definition 3.1]{AR}.

Angular Lipschitz continuity gives, for $|y|>2\max\{|x|,|z|\}$,
\begin{equation}\label{eq:kernel-difference}
 |K(x-y)-K(z-y)|\leq C|x-z|\,|y|^{-n-2s-1}.
\end{equation}
Consequently, if $|U(y)|\leq A(1+|y|)^{1+s+\tau}$ and $\tau<s$, then
\begin{equation}\label{eq:compensated-tail}
 \int_{|y|>R}|U(y)|\,|K(x-y)-K(z-y)|\,dy
       \leq CA|x-z|R^{\tau-s}
\end{equation}
for $x,z$ in a fixed ball and $R$ sufficiently large. The radial integral is
$\int_R^\infty r^{\tau-s-1}\,dr$. In particular, truncations can change
the operator by a divergent constant, but their spatial oscillations converge.

We use the following local estimates from \cite{FRRO}; theorem numbering in
that reference refers to its specified author manuscript. A weak solution
with zero local exterior on a known flat boundary, bounded right-hand side
and bounded weighted tail is locally $C^s$ up to that boundary
\cite[Proposition 2.6.4]{FRRO}. For a continuous distributional solution the
same assertion follows from \cite[Proposition 2.6.13]{FRRO}. In the latter
statement choose its boundary exponent to be $s/2$; a bounded right-hand side
controls the required weighted right-hand side. These results are used on
fixed flat charts, away from any corner of the auxiliary domain.

In the interior, a bounded cutoff solution of $Lw=F$ with $F\in C^\gamma$
is $C^{2s+\gamma}$, provided $2s+\gamma$ is noninteger and the kernel has
the integrated $C^\gamma$ regularity in \cite[Definition 2.1.22]{FRRO}
\cite[Theorem 2.4.1]{FRRO}. Here every $0<\gamma<1$ is allowed by
\eqref{eq:kernel-difference}. The weighted-tail version will always be
applied after the compensation just described.

There is one useful way of controlling a missing constant. Suppose $LU=0$
modulo constants in a fixed half-ball and $U$ has the growth in
\eqref{eq:compensated-tail}. For a fixed compact cutoff $U_0=\chi U$, write
\[
 LU_0=d+G,
\]
where $G$ has bounded Lipschitz oscillation, by \eqref{eq:compensated-tail}.
Normalize $G$ to vanish at one interior point. Testing against a fixed smooth
function $\zeta$ supported strictly inside the half-ball, with
$\int\zeta=1$, gives
\begin{equation}\label{eq:constant-test}
 d=\int_{\R^n} U_0 L\zeta-\int G\zeta.
\end{equation}
Both terms are bounded by $CA$. Indeed $L\zeta\in L^1$, using the
second difference of $\zeta$ near zero and integrability of $K$ at infinity.
Thus the flat $C^s$ estimate applies with a controlled full right-hand side.
This argument also applies to equations whose right-hand-side oscillation
tends to zero.

These observations justify compactness with varying kernels. Uniformly
Lipschitz angular densities have a uniformly convergent subsequence on the
sphere. On finite annuli this convergence passes the operator to the limit;
\eqref{eq:compensated-tail} controls the remaining tails. No convergence of
angular derivatives is needed. Equivalently, one may test against compact
smooth functions of integral zero: their images under $L$ have the extra
far decay $|y|^{-n-2s-1}$.

\subsection{A higher half-space Liouville theorem}

For one-variable functions, polar coordinates and evenness give
\begin{equation}\label{eq:one-dimensional}
 L(F(x\cdot e))=A_e L_sF(x\cdot e),\qquad
 A_e=\frac12\int_{\Sn}a(\theta)|\theta\cdot e|^{2s}\,d\theta>0,
\end{equation}
where
$L_sF(t)=\PV\int_\R(F(t)-F(t+y))|y|^{-1-2s}\,dy$.
The same formula holds for compensated operator differences.
In particular $L(t_+^s)=0$ in $\Hh$
\cite[Lemma 2.6.2 and Lemma B.1.5]{FRRO}.
We will also use
\begin{equation}\label{eq:affine-profile-harmonic}
 L\bigl((a_0+b\cdot x)t_+^s\bigr)=0
                 \quad\hbox{modulo constants in }\Hh.
\end{equation}
For the normal linear term, differentiate $t_+^{1+s}$ spatially and use
$L(t_+^s)=0$. For a tangential index $i<n$, vary the unit normal along
\[
 e(\varepsilon)=\frac{e_n+\varepsilon e_i}{|e_n+\varepsilon e_i|},
 \qquad
 \left.\frac{d}{d\varepsilon}(x\cdot e(\varepsilon))_+^{1+s}
                                      \right|_{\varepsilon=0}
                         =(1+s)x_i t_+^s.
\]
Use radial cutoffs of radius $R$ first and differentiate their spatial
operator oscillation with respect to this normal parameter. The derivative tail
errors on compact sets are $O(R^{-s})$. In the directional calculation the
factor $|\theta\cdot e|^{s-1}$ is integrable on the sphere. Differentiation
of the finite integrals, followed by $R\to\infty$, shows that the spatial
oscillation of the operator vanishes, which is exactly
\eqref{eq:affine-profile-harmonic}.

\begin{lemma}\label{lem:liouville}
Let $0<\tau<\min\{s,1-s\}$. Suppose that $U$ is continuous on $\R^n$,
vanishes in $\{t\leq0\}$, solves $LU=0$ modulo constants in $\Hh$, and
satisfies $|U(x)|\leq A(1+|x|)^{1+s+\tau}$. Then
\[
                  U(x)=(a_0+b\cdot x)t_+^s.
\]
\end{lemma}

\begin{proof}
Apply the cutoff estimate and \eqref{eq:constant-test} to
$R^{-1-s-\tau}U(R\,\cdot)$, and then the flat boundary estimate. This gives
\begin{equation}\label{eq:liouville-Cs}
                 [U]_{C^s(B_R)}\leq CA R^{1+\tau},\qquad R\geq1.
\end{equation}
A tangential increment, divided by $|h|^s$, has growth at most
$1+\tau$, vanishes on the lower half-space, and still solves the zero
equation modulo constants. The same estimate can be applied repeatedly;
the compensated-tail calculation remains valid as the growth decreases.
For a fixed number $j$ with $js>1+\tau$, the $C^s(B_R)$ seminorm of a
$j$-fold tangential increment is bounded by $CR^{1+\tau-js}$.
Sending $R$ to infinity shows that this increment is constant, and its
vanishing on the lower half-space makes it zero.

Thus $U(\cdot,t)$ is a polynomial in $x'$. For completeness, convolve in
$x'$, divide the vanishing finite differences by their step sizes, and let
the steps tend to zero. The resulting smooth function has zero derivatives
of order $j$; passing back through the convolution gives the assertion for
the continuous function. Since $1+s+\tau<2$, its degree in $x'$ is at most
one. Write
\[
              U(x',t)=A_0(t)+\sum_{i=1}^{n-1}x_i B_i(t).
\]
Evaluating this identity at $x'=(1+|t|)e_i$ and at $x'=0$ bounds
$|B_i(t)|\leq CA(1+|t|)^{s+\tau}$. Tangential differences of the equation
show that $LB_i=0$ modulo constants. This operator is now an ordinary
convergent integral, because $s+\tau<2s$, and its value is a constant.
Scaling by $R$ and testing inside the half-space bounds that constant by
$CR^{\tau-s}$, so it is zero. The lower-growth half-space Liouville theorem
\cite[Theorem 2.7.2]{FRRO} gives $B_i(t)=b_i t_+^s$.

Subtract these tangential terms using \eqref{eq:affine-profile-harmonic}.
The remaining one-dimensional $A_0$ is zero for $t\leq0$, has growth
$1+s+\tau$, and solves $L_sA_0=0$ modulo constants by
\eqref{eq:one-dimensional}. The one-dimensional higher Liouville theorem
\cite[Proposition 3.11, $k=1$]{AR} gives
$A_0(t)=(a_0+b_n t)t_+^s$.
\end{proof}

\subsection{The quotient estimate}

\begin{proposition}\label{prop:flat-quotient}
Fix $0<\alpha<\min\{s,1-s\}$, put $\gamma=1-s+\alpha$, and let
$0<\tau<\alpha$. If $q$ is bounded on $\R^n$, continuous in $B_1$, and
\[
 Lq=F\quad\hbox{in }B_1\cap\Hh,\qquad
 q=0\quad\hbox{in }B_1\cap\{t\leq0\},
 \qquad F\in C^\gamma(\overline{B_1\cap\Hh}),
\]
then $q/t^s$ extends to $C^{1,\tau}(\overline{B_{1/2}\cap\Hh})$ and
\begin{equation}\label{eq:flat-quotient}
 \left\|\frac{q}{t^s}\right\|_{C^{1,\tau}(\overline{B_{1/2}\cap\Hh})}
     \leq C\bigl(\|q\|_{L^\infty(\R^n)}+\|F\|_{C^\gamma}\bigr).
\end{equation}
Only local vanishing on the lower half-space is assumed.
\end{proposition}

\begin{proof}
Normalize the right side to one. We first prove, at every flat boundary
point $z$ in a smaller ball, an affine approximation
\begin{equation}\label{eq:boundary-affine}
 \|q-Q_z t_+^s\|_{L^\infty(B_r(z))}\leq Cr^{1+s+\tau},
             \qquad Q_z(x)=a_z+b_z\cdot(x-z),
\end{equation}
with bounded coefficients. Translate $z$ to zero. If the estimate fails,
take a sequence $q_j,L_j$ with the fixed structural bounds. Let
$Q_{j,r}(x)=a_{j,r}+b_{j,r}\cdot x$ be the least-squares affine fit to
$q_j/t^s$ on $B_r\cap\Hh$, that is, the
minimizer of $\int_{B_r\cap\Hh}|q_j-Q t^s|^2$. The Gram matrix of
$\{t^s,x_1t^s,\ldots,x_nt^s\}$ on the unit half-ball is positive definite.
Define
\[
 \Theta(r)=\sup_j\sup_{\rho\geq r}
       \rho^{-1-s-\tau}\|q_j-Q_{j,\rho}t_+^s\|_{L^\infty(B_\rho)}.
\]
The failure gives $\Theta(r)\to\infty$ as $r\downarrow0$. Choose scales
$r_j\downarrow0$ and corresponding solutions for which the normalized
error is at least $\Theta(r_j)/2$, and set
\[
 V_j(x)=\frac{q_j(r_jx)-Q_{j,r_j}(r_jx)r_j^s t_+^s}
                  {r_j^{1+s+\tau}\Theta(r_j)}.
\]
Finite-dimensional norm equivalence at two consecutive scales gives
\[
 |a_{j,r}-a_{j,2r}|\leq C\Theta(r)r^{1+\tau},\qquad
 |b_{j,r}-b_{j,2r}|\leq C\Theta(r)r^\tau.
\]
Telescoping from $r_j$ to $Rr_j$ therefore yields
\begin{equation}\label{eq:Campanato-growth}
 \|V_j\|_{L^\infty(B_R)}\leq CR^{1+s+\tau}\quad(R\geq1),
 \qquad \|V_j\|_{L^\infty(B_1)}\geq\tfrac12.
\end{equation}
When $Rr_j$ exceeds the original ball, boundedness of $q_j$ and the same
coefficient estimates give this bound as well.

The subtracted affine multiple of $t_+^s$ has zero operator oscillation in
$\Hh$ by \eqref{eq:affine-profile-harmonic}. Thus the right-hand-side
oscillation in the equation for $V_j$, on each fixed half-ball, is bounded by
\[
                C_R\frac{r_j^{\alpha-\tau}}{\Theta(r_j)}\longrightarrow0.
\]
The cutoff test \eqref{eq:constant-test} and the flat $C^s$ estimate give
compactness up to the flat boundary. After extracting the angular densities,
\eqref{eq:compensated-tail} passes the equation to a locally uniform limit
$V$. It vanishes on the lower half-space and has the growth
\eqref{eq:Campanato-growth}. Lemma~\ref{lem:liouville} makes it an affine
polynomial times $t_+^s$. But least-squares orthogonality gives
\[
       \int_{B_1\cap\Hh}V_j(x)Q(x)t^s\,dx=0
                 \quad\hbox{for every affine }Q.
\]
Passing to the limit and choosing its coefficient polynomial shows $V=0$,
contradicting \eqref{eq:Campanato-growth}. This proves
\eqref{eq:boundary-affine}; dyadic telescoping gives the limiting
coefficients $a_z,b_z$ and their boundedness.

It remains to recover the quotient's normal derivative. Let $x_0$ have
height $d>0$, let $z=(x_0',0)$, and put $w=q-Q_z t_+^s$.
Estimate~\eqref{eq:boundary-affine}, extended away from $z$ by boundedness
and the coefficient bounds, gives
$|w(x)|\leq C|x-z|^{1+s+\tau}$.
Rescale at $d$ and divide by $d^{1+s+\tau}$. On a fixed interior ball
around $e_n$ the remainder and its compensated tail are bounded, and the
$C^\gamma$ oscillation of its right-hand side is at most
$Cd^{\alpha-\tau}$. After a fixed cutoff, \eqref{eq:constant-test} controls
the remaining scalar. The interior estimate gives a uniform
$C^{1+s+\alpha}$ norm there. Dividing by the smooth positive function
$t^s$ and scaling back gives, with $h=q/t^s$,
\[
 |h(x_0)-Q_z(x_0)|\leq Cd^{1+\tau},\qquad
 |\nabla h(x_0)-b_z|\leq Cd^\tau,
\]
and the corresponding interior $C^\tau$ bound for $\nabla h$.
Comparing the affine approximations on overlapping half-balls gives
$|b_z-b_{\widetilde z}|\leq C|z-\widetilde z|^\tau$ for boundary points
$z,\widetilde z$, and the compatible first-order estimate for
$a_z,a_{\widetilde z}$. The interior and boundary cases together prove
\eqref{eq:flat-quotient}.
\end{proof}

\section{Improvement of flatness for the actual solution}\label{sec:flatness}

Put $p=1+s$ and $P_{\kappa,e}(x)=\kappa(x\cdot e)_+^p$. For a growing
difference from such a profile, the tail that controls the oscillation of
the operator is
\begin{equation}\label{eq:tail-definition}
             \T(w)=\int_{|y|>1}|w(y)|\,|y|^{-n-2s-1}\,dy.
\end{equation}
The extra power in this weight is dictated by
\eqref{eq:kernel-difference}. In particular $\T(P_{\kappa,e})<\infty$,
although its raw operator diverges.

We first recall the precise compactness input. For a bounded nonnegative
solution of $\min\{Lw-f,w\}=0$ in $B_1$, for an exponent
$\vartheta\in(\max\{1,2s\},1+s)$ and $\sigma\in(0,s)$,
\cite[Lemma 6.5]{RTW} gives a local $C^q$ bound for some
$q>\max\{1,2s\}$ in terms of
\begin{equation}\label{eq:rough-compactness}
          [f]_{C^{\vartheta-2s}(B_1)}+
                           \|w\|_{L^\infty_{1+s+\sigma}}.
\end{equation}
The bound uses the forcing seminorm, without its constant value. We apply
this estimate to finite cutoffs and rescalings of actual solutions. The
classification \cite[Theorem 1.5]{RTW} says the following: if a globally
nonnegative, locally Lipschitz and locally $C^{2s+\varepsilon}$ function
$V$ satisfies, for every $R\geq1$ and $h\in\R^n\setminus\{0\}$,
\begin{equation}\label{eq:classification-hypotheses}
 \|\nabla V\|_{L^\infty(B_R)}\leq CR^{s+\sigma},\quad
 L\left(\frac{V(\cdot+h)-V}{|h|}\right)\geq0\ \hbox{in }\{V>0\},
 \quad 0<\sigma<\min\{s,1-s\},
\end{equation}
then it is convex. For our homogeneous kernels, if
$z_0\in\partial\{V>0\}$ and $L(\nabla V)=0$ in its positive phase, then
$V(x)=\kappa((x-z_0)\cdot e)_+^{1+s}$.

\begin{proposition}\label{prop:improvement}
Fix $0<\tau'<\tau<\alpha<\min\{s,1-s\}$ and
$0<\kappa_0<\kappa_1<\infty$. There are $\rho\in(0,1/100)$ and
$\delta_*>0$, depending only on these parameters and the structural bounds,
with the following property. Suppose $v$ is bounded individually on
$\R^n$, nonnegative there, locally $C^q$ for some $q>\max\{1,2s\}$, and
\begin{gather}\label{eq:slab-input}
 \min\{Lv-f,v\}=0\ \hbox{in }B_4,\qquad 0\in\partial\{v>0\},\qquad
 [f]_{C^\gamma(B_4)}\leq\eta\leq\delta\leq\delta_*,\quad \delta>0,
 \\
 \kappa(t-\delta)_+^p\leq v(x)\leq\kappa(t+\delta)_+^p
                   \quad(x\in B_2),\qquad
 \T(v-P_{\kappa,e_n})\leq\delta,\qquad
 \kappa_0\leq\kappa\leq\kappa_1,\notag
\end{gather}
where $\gamma=1-s+\alpha$. Then there are $\kappa',e'$ with
$|\kappa'-\kappa|+|e'-e_n|\leq C\delta$ such that
$v_\rho(x)=\rho^{-p}v(\rho x)$ satisfies
\begin{gather}\label{eq:slab-output}
 \kappa'(x\cdot e'-\delta')_+^p\leq v_\rho(x)
                    \leq\kappa'(x\cdot e'+\delta')_+^p
                         \quad(x\in B_2),\\
 \T(v_\rho-P_{\kappa',e'})\leq\delta',\qquad
             \delta'=\delta\rho^{\tau'} .\notag
\end{gather}
The rescaled forcing seminorm is at most $\eta\rho^\alpha\leq\delta'$.
The constants are independent of the individual global bound of $v$.
\end{proposition}

\subsection{A strip decomposition and the forcing constant}

We prove Proposition~\ref{prop:improvement} by contradiction as
$\delta\downarrow0$. First we obtain estimates valid for each admissible
$v$. Choose $\chi$ supported in $B_2$, equal to one on $B_1$, and put
$P=\kappa t_+^p$ and $W=\chi(v-P)$. In the known positive half-ball
$D_\delta=B_{1/2}\cap\{t>\delta\}$, direct subtraction of the operator at
zero gives
\begin{equation}\label{eq:W-equation}
 LW=D+g,\qquad
 g(x)=f(x)-f(0)+\int(1-\chi(y))(v-P)(y)
                         (K(x-y)-K(-y))\,dy.
\end{equation}
Here $D$ is a scalar. At the actual free-boundary point zero, $Lv(0)=f(0)$
by continuity of the operator from positive points, since $q>2s$.
The compensated difference $LP(x)-LP(0)$ is zero for $t>0$, by
\eqref{eq:affine-profile-harmonic} and the derivative of $P$.
Formula~\eqref{eq:kernel-difference} and the tail assumption imply
\begin{equation}\label{eq:g-bound}
                      \|g\|_{C^\gamma(D_\delta)}\leq C\delta.
\end{equation}
The value bounds also imply, on the support of $\chi$,
\begin{equation}\label{eq:W-values}
 |W|\leq C\delta^p\quad(|t|\leq\delta),\qquad
 |W|\leq C\delta(t+\delta)^s\quad(t\geq\delta),\qquad
 W=0\quad(t\leq-\delta).
\end{equation}

We need a bound for $D$, whose size is not controlled by the forcing
seminorm alone. The harmonic influence of a thin exterior strip provides
it. Exterior data bounded by $M_0\delta^s$ on
$-\delta<t<\delta$ have harmonic extension to $D_\delta$ bounded by
\begin{equation}\label{eq:strip-barrier}
 E_\delta(t)=\frac{\sin(\pi s)}{\pi}M_0\delta^s(t-\delta)^s
       \int_{-\delta}^{\delta}\frac{(\delta-y)^{-s}}{t-y}\,dy,
                         \qquad t>\delta.
\end{equation}
On $t\leq\delta$ this function is the strip datum
$M_0\delta^s\mathbf1_{(-\delta,\delta)}$. The formula is the half-line
Poisson formula. It follows from the one-dimensional interval formula
\cite[Proposition 1.10.10]{FRRO} by sending the right endpoint to infinity:
the ratio
\[
 \frac{(t-\delta)(2R-t+\delta)}{(\delta-y)(2R+\delta-y)}
                       \longrightarrow\frac{t-\delta}{\delta-y}
\]
has an integrable majorant on the strip. The limit equation follows by
bounded stability or direct verification. Formula~\eqref{eq:one-dimensional}
makes this a barrier for the same anisotropic operator. It satisfies
\begin{equation}\label{eq:strip-decay}
 E_\delta\leq M_0\delta^s,
 \qquad E_\delta(t)\leq CM_0\delta\,t^{s-1}\quad(t\geq2\delta).
\end{equation}

Let $\tau_\delta$ be the torsion function of $D_\delta$, with zero exterior
datum and $L\tau_\delta=1$. At $x_\delta=2\delta e_n$,
\begin{equation}\label{eq:torsion}
                  c\delta^s\leq\tau_\delta(x_\delta)\leq C\delta^s.
\end{equation}
For the lower bound compare with the unit-forcing torsion function of a
fixed-radius ball inside $D_\delta$ and tangent to its flat boundary. A fixed
smooth interior bump gives a uniform positive interior lower bound. The
Hopf estimate and its proof in \cite[Proposition 2.6.6]{FRRO} then give the
lower bound in \eqref{eq:torsion}. For the upper bound, a nonnegative smooth
cutoff equal to one on $B_1$ is a supersolution after multiplication by a
structural constant, so $\tau_\delta$ is uniformly bounded. The flat
boundary $C^s$ estimate gives the desired upper bound.

All the auxiliary Dirichlet problems here exist in the weak energy class
by \cite[Theorem 2.2.24]{FRRO}; $n\geq2$ and $s<1$ imply $n>2s$.
The datum $W$ is compactly supported and $C^1$ for each solution. Split its
exterior values by smooth cutoffs into a strip part near zero and a part
separated from the local flat boundary. The partitions have finite-energy
extensions; in particular no artificial jump is made at the curved boundary.
Weak comparison follows by testing the negative part in the symmetric energy
form. The weak boundary estimate \cite[Proposition 2.6.4]{FRRO} gives the
continuous representatives before the continuous-solution boundary theorem
or Proposition~\ref{prop:flat-quotient} is applied.

At $x_\delta$, the harmonic extension of the strip part of $W$ is
$O(\delta^p)$, by \eqref{eq:strip-decay} with $M_0=C\delta$.
The separated harmonic part has norm $O(\delta)$, vanishes on the local flat
exterior, and is therefore also $O(\delta^p)$ there. The zero-exterior
solution of $Lb=g$ is bounded in absolute value by
$\|g\|_\infty\tau_\delta$. Representing $W$ as these three terms plus
$D\tau_\delta$, and using \eqref{eq:W-values}, gives
\[
 |D|\tau_\delta(x_\delta)
      \leq C\delta^p+C\delta\tau_\delta(x_\delta).
\]
It follows that
\begin{equation}\label{eq:D-bound}
                              |D|\leq C\delta.
\end{equation}

\subsection{Outer quotient and inner profile}

Write $z=W/\delta$. Let $e$ be the harmonic extension carrying its strip
datum, with a smooth horizontal cutoff equal to one near zero and supported
away from the curved boundary, and set $q_0=z-e$. It is uniformly bounded,
has zero local exterior on $t\leq\delta$, and satisfies
$Lq_0=D/\delta+g/\delta$. Proposition~\ref{prop:flat-quotient}, after
translation and a fixed rescaling, gives
\[
 q_0(x',t)=(t-\delta)^s H_\delta(x',t-\delta),\qquad
                    \|H_\delta\|_{C^{1,\tau}}\leq C
\]
on a fixed smaller half-ball. The strip bound gives
$|e(x)|\leq C\delta t^{s-1}$ for $t\geq2\delta$. Hence
\begin{equation}\label{eq:outer-expansion}
 v(x)=\kappa t^p+\delta(t-\delta)^sH_\delta(x',t-\delta)
                           +O(\delta^2t^{s-1}),\qquad t\geq2\delta.
\end{equation}

Now let $\delta_j\downarrow0$ and choose arbitrary actual free-boundary
points $x_j=(x_j',\delta_j h_j)$ in a fixed small ball, so $|h_j|\leq1$.
Pass to a subsequence with $x_j'\to x_*'$, $h_j\to h$ and
$\kappa_j\to\kappa$. Center at $x_j$, rescale by $\delta_j$, and divide
values by $\delta_j^p$. For applying \eqref{eq:rough-compactness}, first
multiply the actual $v_j$ by a nonnegative fixed cutoff supported in $B_2$
and equal to one near all these points. Its rescaling agrees with the actual
solution on every fixed ball eventually. The sandwich on the cutoff support
gives a uniform global value bound
\begin{equation}\label{eq:inner-value-growth}
                    |V_j(y)|\leq C(1+|y|)^p.
\end{equation}
The separated-tail forcing has Lipschitz oscillation bounded on a fixed
rescaled ball by $C\delta_j^s$. This uses
$\T(v_j-P_j)\leq\delta_j$ and the finite first-order tail of $P_j$.
The original forcing contributes at most $\eta_j\delta_j^\alpha$ to its
$C^\gamma$ seminorm. Thus \eqref{eq:rough-compactness} gives compactness
in $C^1_{\rm loc}$ and in a lower $C^q_{\rm loc}$ with $q>2s$.
Let $V$ be the nonnegative limit.

For every fixed $R$, dilate the finite $V_j$ by $R$, divide by $R^p$, and
apply \eqref{eq:rough-compactness} before letting $j$ tend to infinity.
The weighted value norm is bounded by \eqref{eq:inner-value-growth}, and
the forcing oscillation tends to zero for that fixed $R$. This gives
\begin{equation}\label{eq:inner-gradient-growth}
                       \|\nabla V\|_{L^\infty(B_R)}\leq CR^s.
\end{equation}
It does not require an ordinary operator on the unbounded profile $V$.
To pass complementarity to increments, let $f_j$ be the forcing of the
uncutoff $v_j$. Whenever $v_j(x)>0$ and both $x,x+\delta_jh$ lie in $B_4$,
the finite equation gives
\[
 L_j\bigl(v_j(\cdot+\delta_jh)-v_j\bigr)(x)
                         \geq f_j(x+\delta_jh)-f_j(x).
\]
There is equality if $v_j(x+\delta_jh)>0$ as well. On each fixed inner
ball the cutoff agrees with the actual rescaling. Its additional forcing
has the vanishing separated-tail oscillation estimated above. After
rescaling and division by $|h|$, the right-hand-side differences tend to
zero for fixed $h\ne0$. For the far part, transfer a signed
increment to a difference of kernels with a smooth cutoff; its error is
$O(|h|R^{-s})$ by the value growth. Passing first to the local limit and
then sending the tail radius $R$ to infinity gives the increment inequality
in \eqref{eq:classification-hypotheses}. On a compact set where both shifted
points are positive, the limiting inequality is an equality. Sending the
shift to zero gives $L(\nabla V)=0$ in the positive phase, with the
viscosity interpretation supplied by local interior regularity.
The classification following \eqref{eq:classification-hypotheses} applies.
The inherited slab bounds at every fixed radius force the resulting profile
to have direction $e_n$ and amplitude $\kappa$.

We must still retain the actual free-boundary point at the origin. A direct
compensated calculation supplies this. For $t<0$,
\[
 L_s(t_+^s)=-\Bfun(s+1,s)(-t)^{-s},\qquad
 \Bfun(s+1,s)=\int_0^\infty z^s(1+z)^{-1-2s}\,dz>0.
\]
Differentiate large cutoffs of $t_+^p$, use their $O(R^{-s})$ derivative
tail error, and integrate from zero to a fixed negative point. Since
$p>2s$, the resulting difference is continuous at zero. The compensated
operator of $P_{\kappa,e}$ at signed distance $-d<0$, relative to its
positive-side value, is
\begin{equation}\label{eq:negative-cusp}
                 A_e\kappa\frac{p\Bfun(s+1,s)}{1-s}d^{1-s}>0.
\end{equation}
Here $A_e>0$ is the one-dimensional factor defined in
\eqref{eq:one-dimensional}.
If the origin were strictly inside the contact region of the limiting
profile, choose positive points $y_j\to0$ of $V_j$ and a fixed positive
point $z$ of $V$. Local $C^q$ convergence with $q>2s$ and compensated tail
convergence give
\[
 L_jV_j(y_j)-L_jV_j(z)
             \longrightarrow (LV)(0)-(LV)(z)
                         \quad\hbox{in the compensated sense}.
\]
The left side is a forcing difference and tends to zero. The right side is
strictly positive by \eqref{eq:negative-cusp}, a contradiction. Thus
\begin{equation}\label{eq:centered-inner-limit}
                              V(y)=\kappa(y_n)_+^p.
\end{equation}
The same calculation excludes actual positive pockets on any fixed compact
set strictly below the limiting plane: a sequence of positive points in
such a set would give the identical contradiction.

\subsection{Signed matching and the improved value bounds}

Extract a locally $C^1$ convergent subsequence of the uniformly
$C^{1,\tau}$ functions $H_{\delta_j}$ in \eqref{eq:outer-expansion}, with
limit $H_\infty$. Return to the arbitrary actual centers
$x_j=(x_j',\delta_jh_j)$, with $x_j'\to x_*'$ and
$h_j=x_{j,n}/\delta_j\to h$. Fix $T>3$ and evaluate at the original point
\[
             x=x_j+\delta_jTe_n,\qquad t=\delta_j(T+h_j).
\]
Divide the outer expansion \eqref{eq:outer-expansion} by the value scale
$\delta_j^p$. Its left side is
$v_j(x_j+\delta_jTe_n)/\delta_j^p\to\kappa T^p$ by
\eqref{eq:centered-inner-limit}. Letting $j\to\infty$ at this fixed $T$ gives
\begin{equation}\label{eq:matching-finite-T}
 \kappa T^p=\kappa(T+h)^p+(T+h-1)^sH_\infty(x_*',0)
                                        +O(T^{s-1}).
\end{equation}
The error constant is uniform in $T>3$. Only after taking $j\to\infty$
with $T$ fixed do we subtract $\kappa T^p$, divide by $T^s$, and send
$T$ to infinity. We obtain
\begin{equation}\label{eq:signed-matching}
                         H_\infty(x_*',0)=-p\kappa h.
\end{equation}
Because the free-boundary sequence was arbitrary, this matching is uniform
on compact horizontal sets. At the original free-boundary point zero it
gives $H_\infty(0,0)=0$. Multiple components are included in this assertion.

Write $H_\infty(x)=b\cdot x+O(|x|^{1+\tau})$. Choose
\begin{equation}\label{eq:new-profile}
 \kappa_j'=\kappa_j+\delta_j b_n,\qquad
 e_j'=\frac{e_n+\delta_j b'/(p\kappa_j)}
                         {|e_n+\delta_j b'/(p\kappa_j)|}.
\end{equation}
The first variation of $P_{\kappa_j',e_j'}$ is
$P_{\kappa_j,e_n}+\delta_j t^s(b\cdot x)$. The quotient Taylor remainder
on $B_{2\rho}$ is bounded by $C\rho^{1+\tau}$. By
\eqref{eq:signed-matching}, every actual free-boundary point lies within
this normalized height error, plus $o(1)$, of the new plane. Choose $\rho$
small enough that $C\rho^{\tau-\tau'}$ is small.

We verify value ordering in two regions. If
$t\geq M_1\delta_j/\rho^{1+\tau'}$, the strip error in
\eqref{eq:outer-expansion}, divided by $\delta_jt^s$, is at most
$C\delta_j/t\leq(C/M_1)\rho^{1+\tau'}$. With $M_1$ fixed sufficiently
large, the quotient remainder and profile Taylor error are absorbed by
the shifts $\delta_j\rho^{1+\tau'}$ in the original variables.
The mean-value formula gives both value inequalities there.

In the remaining strip, choose actual free-boundary points on the vertical
fibers. A top such point exists between heights $-\delta_j$ and
$\delta_j$. More explicitly, for a fixed small $r_0$ within the quotient
chart, set $\Gamma_j=\partial\{v_j>0\}\cap B_{r_0}$. For every fixed
$R<\infty$, the unique centered limit \eqref{eq:centered-inner-limit} gives
\[
 \sup_{z\in\Gamma_j}\ \sup_{y\in B_R}
 \left|\frac{v_j(z+\delta_j y)}{\delta_j^p}
                         -\kappa_j(y_n)_+^p\right|\longrightarrow0.
\]
A violating sequence of centers would have a subsequence with that same
unique limit. Choose $\rho$ first and then the finite inner radius
$R=M_1/\rho^{1+\tau'}$, before sending $j$ to infinity. The resulting
value error is $o(\delta_j^p)$,
smaller than the gap between the two positively shifted profiles, whose
required nonzero gap is at least
$c(\delta_j\rho^{1+\tau'})^p$.

For the zero region in the upper comparison one also needs actual contact,
rather than small values. A positive pocket below the improved matching
band would have a boundary point there: on its vertical fiber the solution
is zero at height $-\delta_j$. This contradicts the arbitrary-point
matching \eqref{eq:signed-matching}. The lower comparison is automatic
where its model vanishes. These facts prove the genuine sandwich on
$B_{2\rho}$ for large $j$, without an assumption of vertical monotonicity.

\subsection{Tail improvement and iteration}

Under the scaling in the proposition,
\begin{equation}\label{eq:tail-scaling}
 \T(v_\rho-P_{\kappa',e'})=
       \rho^s\int_{|y|>\rho}|v(y)-P_{\kappa',e'}(y)|
                                      |y|^{-n-2s-1}\,dy.
\end{equation}
On each fixed annulus $\rho<|y|<r_0$, division by $\delta_j$ and passage
to the limit gives
$t_+^s(H_\infty(y)-b\cdot y)$. The original sandwich bounds the quotient
through the vanishing strip, so dominated convergence applies. Its
contribution to \eqref{eq:tail-scaling} is at most
$C\delta_j\rho^\tau$, since $\tau<s$. Outside a fixed small $r_0$, the
original sandwich, the tail assumption and the $O(\delta_j)$ profile change
give at most $C\delta_j\rho^s$. Further decreasing $\rho$ ensures
\[
                       C(\rho^\tau+\rho^s)<\tfrac12\rho^{\tau'}.
\]
Thus the tail in \eqref{eq:slab-output} is at most
$\delta_j\rho^{\tau'}$ for large $j$. Together with the value bounds this
contradicts a sequence violating Proposition~\ref{prop:improvement} and
proves the existence of $\delta_*$.

For iteration, fix $0<\varepsilon<(\kappa_1-\kappa_0)/2$ and start in
the nonempty interval $[\kappa_0+\varepsilon,\kappa_1-\varepsilon]$,
allowing the larger interval $[\kappa_0,\kappa_1]$. Decrease $\delta_*$
so that $C\delta_* /(1-\rho^{\tau'})<\varepsilon$. Then
$\delta_j=\delta_0\rho^{j\tau'}$, and the changes of amplitude and direction
are summable. The value sandwiches imply
\begin{equation}\label{eq:endpoint-value}
                           \|v\|_{L^\infty(B_r)}\leq Cr^p,
                                  \qquad 0<r<1.
\end{equation}
To obtain the gradient, cut off the actual $v$ in $B_2$. Its separated-tail
forcing has uniformly Lipschitz oscillation by the original sandwich and
$\T$, although its constant value need not be small. At scale $r$ divide
values by $r^p$. Estimate~\eqref{eq:endpoint-value} bounds the global
weighted value norm of this cutoff. Its forcing seminorm is bounded by
$C(\eta r^\alpha+r^s)$. Applying \eqref{eq:rough-compactness} to each finite
rescaling gives
\begin{equation}\label{eq:endpoint-gradient}
                           \|\nabla v\|_{L^\infty(B_r)}\leq Cr^s,
                                  \qquad 0<r<1.
\end{equation}

\section{The interior estimate}\label{sec:interior}

We now return to the original obstacle problem and prove
Theorem~\ref{thm:main}. Normalize
\begin{equation}\label{eq:normalization}
             \|\phi\|_{C^\beta(B_1)}+\|u\|_{L^\infty(\R^n)}=1.
\end{equation}
If the left side is zero, the result is immediate. Choose, once and for all,
\begin{equation}\label{eq:exponents}
 \begin{gathered}
 0<\alpha<\min\{\beta-1-s,s,1-s\},\qquad
 b=1+s+\alpha<\min\{\beta,1+2s,2\},\\
 \gamma=1-s+\alpha,\qquad 0<\sigma<\alpha,\qquad
 0<\tau'<\tau<\alpha.
 \end{gathered}
\end{equation}
All these choices depend only on $s,\beta$.

\subsection{Localization of the rough obstacle}

The subcritical part of \cite[Corollary 1.3]{RTW} gives a uniform
$C^q$ bound in $B_{15/16}$ for a fixed
$\max\{1,2s\}<q<1+s$. Only that subcritical statement is used here.
Take $\xi\in C_c^\infty(B_1)$ equal to one on $B_{15/16}$ and set
$\widetilde\phi=\xi\phi$, extended by zero. Its global $C^b$ norm is
uniformly bounded. With $0\leq\chi\leq1$ supported in $B_{7/8}$ and equal
to one on $B_{3/4}$, put
\begin{equation}\label{eq:localized-v}
                         v=\chi(u-\widetilde\phi).
\end{equation}
It is globally nonnegative, bounded and Lipschitz, with uniform bounds.
In $B_{2/3}$ it agrees with $u-\phi$ and solves
\begin{equation}\label{eq:localized-equation}
 \min\{Lv-f,v\}=0,\qquad
 f=-L\widetilde\phi-L((1-\chi)(u-\widetilde\phi)),\qquad
                            \|f\|_{C^\gamma(B_{2/3})}\leq C.
\end{equation}
The mapping estimate \cite[Lemma 2.2.6]{FRRO} gives the first term's
$C^\gamma$ bound, since $b=2s+\gamma$. The second term is a separated tail;
differentiate its kernel once to obtain a Lipschitz bound, with no
derivative of the exterior data. The local exponent $q>2s$ makes the
equation hold both classically and distributionally. At every zero of $v$,
$\nabla v=0$.

\subsection{Coarse entry at a nondegenerate scale}

We record the compactness consequence needed for the scale selection.
Given $c_0>0$ and $\delta_0>0$, there is $\eta_0>0$ such that a bounded,
globally nonnegative solution with
\begin{gather}\label{eq:coarse-hypotheses}
 \min\{Lv-f,v\}=0\ \hbox{in }B_{1/\eta},\quad
 0\in\partial\{v>0\},\quad [f]_{C^\gamma}\leq\eta\leq\eta_0,\\
 \|\nabla v\|_{L^\infty(B_R)}\leq R^{s+\sigma}\quad(R\geq1),\qquad
                  \|\nabla v\|_{L^\infty(B_1)}\geq c_0\notag
\end{gather}
has a profile $P_{\kappa,e}$ with amplitude bounded above and bounded away
from zero, and satisfies
\begin{equation}\label{eq:coarse-sandwich}
 \kappa(x\cdot e-\delta_0)_+^p\leq v(x)
                    \leq\kappa(x\cdot e+\delta_0)_+^p\quad(x\in B_2),
 \qquad \T(v-P_{\kappa,e})\leq\delta_0.
\end{equation}
The amplitude bounds depend only on $c_0$ and the structural parameters.

Indeed, otherwise let $\eta_j\downarrow0$. The gradient growth in
\eqref{eq:coarse-hypotheses} gives weighted value growth of degree
$p+\sigma$. The compactness estimate \eqref{eq:rough-compactness} gives a
locally $C^1$ limit with a local exponent greater than $2s$.
Complementarity passes to the signed increment inequalities as in the
inner argument above. The classification
\eqref{eq:classification-hypotheses} makes the limit a nonzero half-space
profile. The compensated cusp argument \eqref{eq:negative-cusp} retains
the origin as its actual free-boundary point and excludes positive points
in negative compact sets. Thus it is $\kappa(x\cdot e)_+^p$, with a fixed
positive lower amplitude from $c_0$.

In particular $v_j=0$ on $B_2\cap\{x\cdot e\leq-\delta_0/4\}$ for
large $j$. Uniform value convergence, with error $o(\delta_0^p)$, and
the positive lower amplitude give the shifted comparisons
\eqref{eq:coarse-sandwich}. Where the lower profile vanishes its inequality
is automatic. In the remaining region its positive shift gap, and likewise
the upper shift gap above the actual zero region, dominate that error.
For the tail, convergence on fixed annuli is local convergence, while the
part beyond $R$ is bounded by $CR^{\sigma-s}$. Since $\sigma<s$, first
choose $R$ large, then $j$ large. This proves the coarse assertion.

Choose $\delta_0$ below the threshold of Proposition~\ref{prop:improvement}
for comfortable nested amplitude intervals, and then $\eta_0\leq\delta_0$
accordingly. Rotating $e$ to $e_n$ preserves all structural bounds.
The iteration leading to \eqref{eq:endpoint-gradient} gives
\begin{equation}\label{eq:coarse-to-endpoint}
       \|\nabla v\|_{L^\infty(B_r)}\leq Cr^s\qquad(0<r<1)
\end{equation}
whenever \eqref{eq:coarse-hypotheses} holds.

\subsection{First failure of subcritical decay}

Fix an arbitrary actual free-boundary point $z\in B_{5/8}$ and translate
it to zero. There is a fixed $R_{\rm loc}>0$ on which
\eqref{eq:localized-equation} holds, uniformly in $z$. Let $C_1$ bound
the global Lipschitz norm of $v$ and $C_f$ its local forcing seminorm.
Choose $r_*\leq R_{\rm loc}\eta_0/2$ and then $A$ sufficiently large that
\begin{equation}\label{eq:first-scale-choice}
 \frac{C_1}{Ar_*^{s+\sigma}}\leq1,\qquad
                     \frac{C_f r_*^{\alpha-\sigma}}A\leq\eta_0.
\end{equation}
For $r_k=r_*2^{-k}$ define
\begin{equation}\label{eq:dyadic-normalization}
 w_k(x)=\frac{v(r_kx)}{Ar_k^{1+s+\sigma}},\qquad
 f_k(x)=\frac{r_k^{s-1-\sigma}}A f(r_kx).
\end{equation}
Their equations hold in $B_{1/\eta_0}$ and
\[
                 [f_k]_{C^\gamma}\leq C_f r_k^{\alpha-\sigma}/A
                                             \leq\eta_0.
\]
The strict inequality $\alpha>\sigma$ is used here. Let $G_k$ denote
\begin{equation}\label{eq:Gk}
          \|\nabla w_k\|_{L^\infty(B_R)}\leq R^{s+\sigma}
                                      \quad\hbox{for every }R\geq1.
\end{equation}
By \eqref{eq:first-scale-choice}, $G_0$ holds. If every $G_k$ holds, scaling
back gives $\|\nabla v\|_{L^\infty(B_r)}\leq Cr^{s+\sigma}$ and hence
the endpoint estimate.

Otherwise let $K+1$ be the first failing index. The identity
\[
             \nabla w_{K+1}(x)=2^{s+\sigma}\nabla w_K(x/2)
\]
and $G_K$ already give $G_{K+1}$ for $R\geq2$. If
$\|\nabla w_K\|_{L^\infty(B_1)}\leq2^{-s-\sigma}$, they give it also
for $1\leq R\leq2$. Thus
\[
              \|\nabla w_K\|_{L^\infty(B_1)}>c_0=2^{-s-\sigma}.
\]
Apply the coarse assertion with this fixed $c_0$, followed by
\eqref{eq:coarse-to-endpoint}. For $r<r_K$ it yields
\[
 \|\nabla v\|_{L^\infty(B_r)}\leq AC r_K^\sigma r^s
                                           \leq AC r_*^\sigma r^s.
\]
For $r_K\leq r\leq r_*$, $G_K$ gives
$\|\nabla v\|_{L^\infty(B_r)}\leq Ar^{s+\sigma}
\leq Ar_*^\sigma r^s$. Larger fixed radii are controlled by the original
Lipschitz bound. We have therefore proved, uniformly at every actual
free-boundary point $z\in B_{5/8}$,
\begin{equation}\label{eq:one-point-original}
       \|\nabla(u-\phi)\|_{L^\infty(B_r(z))}\leq Cr^s
\end{equation}
for all sufficiently small fixed radii. No nondegeneracy hypothesis remains.

\subsection{From the one-point estimate to the full norm}

Set $E=\{u=\phi\}\cap B_{2/3}$. Take a smooth cutoff $\theta$ supported
in $B_{3/4}$ and equal to one on $B_{2/3}$, and let $U=\theta u$.
The subcritical estimate bounds $W_i=\partial_iU$ globally. On the actual
positive phase inside $B_{2/3}$,
\begin{equation}\label{eq:derivative-equation}
 LU=F_{\rm tail}=-L((1-\theta)u),\qquad
                         LW_i=\partial_iF_{\rm tail}
\end{equation}
in distributions. The separated tail is uniformly Lipschitz by one kernel
derivative, so the last right-hand side is bounded. We have differentiated
neither $L\phi$ nor the exterior data.

At a contact point $z$, $W_i(z)=\partial_i\phi(z)$. At a free-boundary
point, \eqref{eq:one-point-original} and
$\nabla\phi\in C^{s+\alpha}$ give
\begin{equation}\label{eq:derivative-at-contact}
          |W_i(y)-\partial_i\phi(z)|\leq C|y-z|^s
\end{equation}
near $z$. For $x,y\in B_{1/2}$ put $h=|x-y|$. If $x$ is in the positive
phase and $h$ is comparable to or greater than $d=\dist(x,E)$, a closest
contact point and \eqref{eq:derivative-at-contact} give the $Ch^s$ bound
for small $d,h$. For distances bounded below the global derivative bound
suffices.

Suppose next $h$ is much smaller than $d$ and $d$ is small. Choose a
closest contact point $z$, which is an actual interior free-boundary point,
and put
\[
                  Z_i(X)=\frac{W_i(z+dX)-\partial_i\phi(z)}{d^s}.
\]
Estimate~\eqref{eq:derivative-at-contact} near $z$ and boundedness away
from $z$ imply the uniform weighted value growth
$|Z_i(X)|\leq C(1+|X|)^s$. On a fixed interior ball around $(x-z)/d$,
its equation has bounded right-hand side $d^s\partial_iF_{\rm tail}$.
Choose a finite integrability exponent $p_0$ so that
\[
                   s<\mu=2s-n/p_0<\min\{2s,1\}.
\]
The interior estimate \cite[Theorem 2.4.3]{FRRO} gives a uniform $C^\mu$
bound there; the tail growth $s<2s$ is admissible. Consequently
\[
       |W_i(x)-W_i(y)|\leq Cd^s(h/d)^\mu\leq Ch^s.
\]
If $d$ is bounded below, $u$ is harmonic on a fixed ball and the interior
estimate \cite[Theorem 2.4.1]{FRRO}, with the regular kernel tail, gives
the same conclusion. If both points lie in $E$, use
$\nabla u=\nabla\phi$. For a mixed pair with a point in the interior of
$E$, the segment to the positive point meets an actual free-boundary point
within distance $h$ of both. Apply \eqref{eq:one-point-original} there
and then the obstacle-gradient modulus. If the coincidence set is locally
empty, the harmonic case applies.

These cases cover $B_{1/2}$ and show $[\nabla u]_{C^s(B_{1/2})}\leq C$.
The supremum norms of $u$ and $\nabla u$ were already bounded by the
normalization and the subcritical estimate. Restoring
\eqref{eq:normalization} proves
\[
 \|u\|_{C^{1+s}(B_{1/2})}
       \leq C\bigl(\|\phi\|_{C^\beta(B_1)}+\|u\|_{L^\infty(\R^n)}\bigr),
\]
where $C$ depends only on $n,s,\beta,\lambda,\Lambda$ and
$\|a\|_{\Lip(\Sn)}$, as asserted.\hfill$\square$

\section*{AI assistance}
An AI assistant was used to develop the argument, check sources, and draft
this manuscript. AI feedback was also used during revision.
The paper is submitted as a proposed solution for independent
mathematical review.

\begin{thebibliography}{9}
\bibitem{AR}
N. Abatangelo and X. Ros-Oton,
\emph{Obstacle problems for integro-differential operators: higher regularity
of free boundaries}, author manuscript,
\href{https://arxiv.org/abs/1909.10339v2}{arXiv:1909.10339v2} (2019).

\bibitem{CDS}
L. A. Caffarelli, D. De Silva and O. Savin,
\emph{The two membranes problem for different operators},
Ann. Inst. H. Poincar\'e C Anal. Non Lin\'eaire \textbf{34} (2017), 899--932.
\href{https://doi.org/10.1016/j.anihpc.2016.05.006}{doi:10.1016/j.anihpc.2016.05.006}.

\bibitem{FRRO}
X. Fern\'andez-Real and X. Ros-Oton,
\emph{Integro-Differential Elliptic Equations}, Progress in Mathematics 350,
Birkh\"auser, 2024.
The numbered statements cited here are in the author manuscript
\href{https://arxiv.org/abs/2411.12455v1}{arXiv:2411.12455v1}.

\bibitem{RTW}
X. Ros-Oton, C. Torres-Latorre and M. Weidner,
\emph{Semiconvexity estimates for nonlinear integro-differential equations},
Comm. Pure Appl. Math. \textbf{78} (2025),
\href{https://doi.org/10.1002/cpa.22237}{doi:10.1002/cpa.22237}.
The numbered statements cited here refer to
\href{https://arxiv.org/abs/2306.16751v2}{arXiv:2306.16751v2}.
\end{thebibliography}
\end{document}
